\section{问题重述}

\subsection{问题背景}

单波束测深利用声波在海水中的匀速直线传播测量垂直水深，其测深数据沿航迹密集、测线间缺失。多波束测深系统在与航迹垂直的平面内一次发射多个波束，在海底平坦海域能形成以测线为轴、具有一定宽度的全覆盖水深条带。条带的覆盖宽度 $W$ 随换能器开角 $\theta$ 和水深 $D$ 变化；相邻平行条带的重叠率定义为 $\eta=1-d/W$，其中 $d$ 为相邻测线间距，$\eta<0$ 表示漏测。为保证数据完整，相邻条带通常要求 $10\%\sim20\%$ 的重叠率。真实海底起伏较大，若按平均水深或最浅处水深固定布设测线，会在浅水处漏测或深水处冗余。本文研究对象与任务均来自竞赛题面\cite{prob}。

\subsection{问题内容}

\textbf{问题一：}在与测线方向垂直的平面内，海底坡面与水平面夹角为 $\alpha$。建立多波束覆盖宽度及相邻条带重叠率的数学模型；当开角为 $120^\circ$、坡度为 $1.5^\circ$、中心水深为 $70\text{ m}$ 时，计算表 1 所列位置的指标并保存至 \texttt{result1.xlsx}。

\textbf{问题二：}考虑矩形待测海域，测线方向与海底坡面法向在水平面投影的夹角为 $\beta$，建立覆盖宽度模型；当开角为 $120^\circ$、坡度为 $1.5^\circ$、中心水深为 $120\text{ m}$ 时，计算表 2 所列位置与方向下的覆盖宽度并保存至 \texttt{result2.xlsx}。

\textbf{问题三：}在南北长 $2$ 海里、东西宽 $4$ 海里的矩形海域内，中心水深 $110\text{ m}$、西深东浅、坡度 $1.5^\circ$、开角 $120^\circ$。设计一组总长度最短、完全覆盖且相邻重叠率满足 $10\%\sim20\%$ 的测线。

\textbf{问题四：}附件 \texttt{附件.xlsx} 为某海域（南北 $5$ 海里、东西 $4$ 海里）的单波束测深数据。设计多波束测线，使条带尽可能覆盖整个海域、相邻重叠率尽量不超过 $20\%$、测线总长度尽可能短，并计算测线总长度、漏测面积占比与重叠率超过 $20\%$ 部分的总长度。
